% Einheit 2 von 3 – Bayes (bank/bedingte-wahrscheinlichkeit-und-bayes/e2.jsonl, zusammenbau v0.1)
%% TODO Zeitmarke Einheit 2: Marken-Zeile nicht lesbar
%% TODO Zweigzeile Teil 1: was hier gelernt wird (Satz in Schülersprache aus dem Einheitstitel) – Einheit 2
\einheitenkopf[e2]{Einheit 2 von 3 · Bayes}
\zweigzeile{Abitur GK}
\setcounter{aufgabe}{9}

%% TODO Titel: Ich-kann-Satz fehlt in der Bank (Platzhalter: Kettenname) – Nr. 10
%% TODO Anweisung: ein Satz über der ersten Teilaufgabe fehlt in der Bank – Nr. 10
\begin{aufgabe}{Bayes}
%% TODO \rechenplatz{n}: Schrittzahl des Grundfalls fehlt in der Bank (nur bei fester Schreibform, 2.3 b)
\begin{teile}
\teil Ein Baum hat als erste Stufe das Werk (1 oder 2), als zweite Stufe „defekt“ oder „nicht defekt“. Gefragt ist: Wie wahrscheinlich ist ein Gerät aus Werk 1 defekt? Geht das mit dem Baum oder gegen den Baum? \\ \kreuz{mit dem Baum} \\ \kreuz{gegen den Baum}
\teil Das Baumdiagramm beschreibt eine Untersuchung mit einem Test. Eine Person hat ein positives Ergebnis. Mit welcher Wahrscheinlichkeit ist sie krank?

\baumzwei{krank/0.1,gesund/0.9}{positiv/0.8,negativ/0.2}{positiv/0.1,negativ/0.9}
\leerfeld
\teil In einem Lager stammen 30\,\% der Geräte aus Werk 1, der Rest aus Werk 2; aus Werk 1 sind 4\,\% defekt, aus Werk 2 sind es 2\,\%. Gib an, welche Wahrscheinlichkeit der Term $\frac{0{,}3 \cdot 0{,}04}{0{,}3 \cdot 0{,}04 + 0{,}7 \cdot 0{,}02}$ beschreibt.
\teil Ein Onlineshop verschickt 30\,\% seiner Pakete mit dem Dienst X. 6\,\% aller Pakete werden mit X verschickt und kommen beschädigt an, 7\,\% aller Pakete werden anders verschickt und kommen beschädigt an. Ein Paket kommt beschädigt an. Ist die Wahrscheinlichkeit, dass es mit X verschickt wurde, größer als 40\,\%?
\teil Für eine heikle Frage würfelt jede befragte Person verdeckt. Bei 1 bis 4 beantwortet sie ehrlich „Hast du schon einmal abgeschrieben?“, bei 5 oder 6 ehrlich „Hast du noch nie abgeschrieben?“. 30\,\% der Befragten haben schon einmal abgeschrieben. Eine Person antwortet mit Ja. Mit welcher Wahrscheinlichkeit hat sie schon einmal abgeschrieben? \leerfeld
\teil In Beutel A sind doppelt so viele rote wie blaue Kugeln, in Beutel B 10 rote und 5 blaue, in Beutel C 2 blaue und r rote. Ein Beutel wird zufällig gewählt, jeder mit $\frac{1}{3}$, und daraus eine Kugel gezogen; sie ist blau. Weise nach, dass sie mit der Wahrscheinlichkeit $\frac{\frac{1}{3} \cdot \frac{2}{r+2}}{\frac{1}{3} \cdot \frac{2}{r+2} + \frac{2}{3} \cdot \frac{1}{3}}$ aus Beutel C stammt, und bestimme r, wenn diese Wahrscheinlichkeit $\frac{1}{4}$ beträgt. \hfill (Abitur 2024 LK)
\end{teile}
\end{aufgabe}

%% TODO Titel: Ich-kann-Satz fehlt in der Bank (Platzhalter: Kettenname) – Nr. 11
%% TODO Anweisung: ein Satz über der ersten Teilaufgabe fehlt in der Bank – Nr. 11
\begin{aufgabe}{Bayes – Fehler finden, Begründen, Darstellungswechsel, Anwendung}
\begin{teile}
\teil 4\,\% der Menschen haben eine bestimmte Krankheit. Ein Test ist bei Kranken mit 95\,\% positiv, bei Gesunden mit 5\,\%. Paul sagt: „Wer positiv getestet ist, ist mit 95\,\% krank.“ Finde den Fehler.
\teil Begründe, warum im Nenner des Bayes-Quotienten die Summe aller Pfade zum eingetretenen Ereignis steht.
\teil K bedeutet „krank“, T bedeutet „Test positiv“. Übertrage die Wahrscheinlichkeiten aus dem Baumdiagramm in eine Vierfeldertafel.

\baumzwei{K/0.05,nicht K/0.95}{T/0.9,nicht T/0.1}{T/0.1,nicht T/0.9}
\vierfeldertafel{K}{T}{}
\teil Bei einer Vorsorgeuntersuchung haben 1\,\% der untersuchten Frauen eine bestimmte Erkrankung. Der Befund ist bei Erkrankten mit 90\,\% positiv, bei Gesunden mit 9\,\%. Eine Frau erhält einen positiven Befund. Wie wahrscheinlich ist sie tatsächlich erkrankt, und was folgt daraus für das Arztgespräch?
\end{teile}
\end{aufgabe}
